\documentclass[11pt]{article}
\usepackage[margin=1in]{geometry}
\usepackage{booktabs}
\usepackage{longtable}
\usepackage{array}
\usepackage{amsmath}
\usepackage{amssymb}
\usepackage{xcolor}
\usepackage{hyperref}
\usepackage{titlesec}
\usepackage{enumitem}
\usepackage{fancyhdr}
\hypersetup{colorlinks=true, linkcolor=black, urlcolor=blue}
\titleformat{\section}{\large\bfseries}{\thesection}{0.6em}{}
\titleformat{\subsection}{\normalsize\bfseries}{\thesubsection}{0.6em}{}
\setlength{\parskip}{0.5em}
\setlength{\parindent}{0pt}
\pagestyle{fancy}\fancyhf{}
\rhead{\small Alter multitensor GSVD --- independent replication}
\lhead{\small Prepared for R. Stevens}
\cfoot{\thepage}

\newcommand{\C}{\textit{C}}

\title{\textbf{Independent Replication of the Alter \emph{et al.} \\
Multitensor GSVD Predictor for Neuroblastoma Prognosis}\\[4pt]
\large A method-merit and generalizability study on open data}
\author{Prepared for Rick Stevens, Argonne National Laboratory}
\date{2026-06-23}

\begin{document}
\maketitle
\thispagestyle{fancy}

\begin{abstract}
\noindent
We report an independent, end-to-end replication study of the multitensor
generalized singular value decomposition (GSVD / HO-GSVD) method of Alter
\emph{et al.} (\emph{APL Quantum}, 2026), which derives two orthogonal
tumor-specific copy-number (CN) predictors --- the antisymmetric arraylets
$u_{1,1}$ and $u_{1,101}$ --- from patient-matched tumor-genome and blood-genome
whole-genome profiles in neuroblastoma (NBL), and reports a combined survival
predictor with concordance $\C\approx0.77$--$0.80$ that exceeds and is independent
of every standard-of-care indicator. Because the paper's discovery and validation
data are whole-genome sequencing read-count bins under dbGaP controlled access
(phs000467, TARGET-NBL), we asked, \emph{before} any data exchange with the
authors: does the method's central mathematical structure reproduce on
independent data with an independent implementation, and does its prognostic
claim generalize? Replication of this kind is the ordinary, healthy mechanism by
which a striking single-cohort finding earns durable trust; that is the spirit of
this study. We built a from-scratch GSVD engine, reconstructed the paper's exact
two-genome geometry from open data, and tested it across eight cancer cohorts
chosen for known two-mechanism biology. \textbf{Result, in brief:} the
decomposition and its two-orthogonal-mechanism structure reproduce robustly and
recover real, cancer-specific CN biology --- a genuine and durable methodological
contribution --- while the strong prognostic claim did not generalize to any of
the other seven CN-driven cancers, where the predictor performed at chance and was
outperformed by patient age and an ordinary supervised model. The most
parsimonious reading is that the prognostic result reflects neuroblastoma-specific
copy-number biology and/or in-sample optimism in a single discovery cohort --- both
of which call for the standard remedy: replication on the original data with
cross-validation and a supervised baseline.
\end{abstract}

\section{Background and the replication question}

Alter \emph{et al.} apply a tensor generalization of the GSVD to three
patient-matched layers in neuroblastoma: the tumor genome ($Z_1$), the blood
(germline) genome ($Z_2$), and the tumor transcriptome ($Z_3$), with
$Z_1 = 2{,}831{,}960$ and $Z_2 = 2{,}831{,}959$ one-kilobase \texttt{hg19}
copy-number bins. From the mathematically most ``tumor-exclusive'' triplet they
extract an arraylet $u_{1,1}$ and from the most ``blood-exclusive'' triplet a
second arraylet $u_{1,101}$; these are reported to be (i) mutually orthogonal,
(ii) each survival-predictive, and (iii) jointly more accurate than and
statistically independent of MYCN amplification, INSS stage, age, and COG risk.
The headline numbers are a combined concordance of $\C\approx0.77$ (398-patient
set) up to $0.80$ (range across representations), versus MYCN at
$\C\approx0.70$--$0.73$.

The result is striking, and the data layer that carries it --- tumor-versus-blood
whole-genome read-count bins --- is controlled-access. The scientifically natural
next step is independent replication. We framed two questions:

\begin{enumerate}[leftmargin=1.4em]
  \item \textbf{Method merit.} Does the GSVD decomposition reproduce the
  claimed structure (the two orthogonal tumor-specific arraylets) when
  re-implemented from scratch and applied to the same \emph{type} of data?
  \item \textbf{Generalizability.} Does the prognostic claim
  ($\C\approx0.8$, beating standard of care) hold on independent cohorts of the
  method's native data type?
\end{enumerate}

\section{Everything we tried (chronological)}

This section documents every attempt, including the ones that did not work, since
the path itself is part of the scientific record.

\subsection{Attempt 1 --- Open multiomic proxy on the original disease (Path B)}
\textbf{What:} test the method on the paper's actual disease (neuroblastoma) using
the only open layers available --- TARGET-NBL tumor RNA-Seq $\times$ tumor
methylation (Illumina 450K), patient-matched, $\sim$132 patients, top-5000
variable features each.\\
\textbf{Why:} the WGS CN layer is gated; RNA $\times$ methylation was the closest
open patient-matched two-layer structure.\\
\textbf{Result:} \emph{did not reproduce.} The GSVD arraylets were at chance for
survival ($\C\approx0.45$--$0.50$), the two arraylets were \emph{not} orthogonal
($\cos\approx0.33$), and a plain penalized Cox model on expression principal
components beat everything ($\C=0.738$).\\
\textbf{Lesson:} this was the \emph{wrong data structure}. The paper's signal is
tumor-genome versus blood-genome \emph{copy number}; RNA $\times$ methylation is
two different assay types, not two genomes. A chance result here does not test the
method on its own terms. This motivated finding the correct geometry on open data.

\subsection{Attempt 2 --- Characterizing the controlled-access barrier}
\textbf{What:} determine precisely what is gated and whether any open back-door to
the NBL copy-number signal exists.\\
\textbf{Findings (live GDC API queries):}
\begin{itemize}[leftmargin=1.4em,nosep]
  \item phs000467 $=$ TARGET Neuroblastoma (v23.p8); the $\sim$2.8M-bin CN profiles
  derive from the WGS layer, which is controlled.
  \item TARGET-NBL in the harmonized GDC: \textbf{zero} open copy-number files,
  and in fact \textbf{zero} copy-number files at any access tier --- the NBL CN
  data live in the dbGaP raw layer / legacy TARGET matrix, not in harmonized GDC.
  \item The open TARGET-NBL files ($\sim$1{,}568) are RNA-Seq and methylation only
  --- the same wrong-structure layers Attempt~1 used.
\end{itemize}
\textbf{Lesson:} an exact, bin-identical NBL reproduction is impossible from
public data. It requires either the authors' processed Datasets~1--3 (offered on
request) or a direct dbGaP Data Access Request for phs000467 (for which an
NIH-funded PI is independently eligible). This is a data-availability constraint,
not a verdict on the method.

\subsection{Attempt 3 --- Reconstructing the exact two-genome geometry on open data}
\textbf{Key realization:} the method's prior validation domain (cited in the paper)
was ``tumor and blood whole genomes of 85 astrocytoma'' patients. Brain tumors and
many other cancers \emph{do} expose matched tumor and blood-derived-normal
genome-wide copy number in the \emph{open} GDC tier, via the ``Masked Copy Number
Segment'' (DNAcopy, SNP6) product. Critically, GDC runs DNAcopy \emph{per
aliquot}, so the blood-derived-normal aliquot yields its own independent
genome-wide CN segment file --- a genuine analog of the paper's blood-genome layer
$Z_2$, not a substitute omic.\\
\textbf{What we built:} an independent GSVD/HO-GSVD engine
(\texttt{gsvd\_reference.py}; 10/10 unit tests; reconstruction error $<10^{-8}$),
a survival module (KM, log-rank, Cox, Harrell's \C), and an open-data pipeline
that tiles each segment file into genome-wide bins, autosomal-median-centers per
patient (the paper's convention), and builds
$Z_1=$ tumor-CN-bins $\times$ patients and $Z_2=$ blood-CN-bins $\times$ the same
patients. This reconstructs the paper's exact two-genome geometry on open data.

\subsection{Attempt 4 --- First faithful run: glioma (GBM+LGG), the method's prior domain}
\textbf{What:} 976 patients with matched tumor + blood-normal genome-wide CN
($\sim$11$\times$ the original 85-astrocytoma set).\\
\textbf{Result:} the geometry reproduced --- the two extreme arraylets came out
orthogonal ($\cos=0.0023$) --- but the tumor-exclusive arraylet was at chance for
survival ($\C=0.51$), while age ($\C=0.744$) and a supervised Cox model
($\C=0.737$) clearly beat it.\\
\textbf{Interpretation:} consistent with biology --- adult glioma prognosis is
driven by point mutations and methylation (IDH1/2, MGMT), not copy number, so a
CN-only arraylet should not predict survival there.

\subsection{Attempt 5 --- Direct test of the thresholding step}
\textbf{What:} the paper splits patients at one standard deviation from the
arraylet mean before the survival test. We applied that exact $\sigma$-cutoff to
our data.\\
\textbf{Result:} the cutoff did \emph{not} lift concordance off chance at any
threshold (it stayed $\C=0.50$--$0.51$; at $+1\sigma$ only a single patient fell
in the high group, because the arraylet score had negligible spread).\\
\textbf{Lesson:} the discovery-cohort number is not explained by a thresholding
artifact we could reproduce --- which points \emph{away} from any methodological
sleight and toward the two ordinary explanations in \S\ref{sec:interp}.

\subsection{Attempt 6 --- The two-mechanism panel: eight cancers}
\textbf{What:} the paper's distinctive claim is not merely ``CN predicts
survival'' but that the method recovers \emph{two orthogonal, independent
mechanisms}. To test that, we ran seven additional cohorts, each selected because
it has two \emph{named, biologically-distinct} CN programs to map the arraylets
onto. Each cohort was run through the identical pipeline; for each we recorded both
arraylet survival concordances, their combined concordance, the cross-arraylet
orthogonality, age and supervised-Cox baselines, and the top-loading genomic
regions of each arraylet.

\begin{center}
\small
\begin{tabular}{@{}lll@{}}
\toprule
\textbf{Cohort} & \textbf{Matched pairs} & \textbf{Two known orthogonal axes} \\
\midrule
Glioma (GBM+LGG) & 976 & astrocytoma (the method's prior domain) \\
Breast (BRCA)    & 978 & amplifier (8q/MYC, 11q13/CCND1, 17q/HER2) vs 1q-gain/16q-loss \\
Ovarian (OV)     & 441 & HRD/BRCA-scar vs CCNE1 amplification (mutually exclusive) \\
Sarcoma (SARC)   & 232 & whole-genome-doubling vs focal 12q (MDM2/CDK4) \\
Endometrial (UCEC)& 497 & CN-high serous-like vs CN-low/MSI \\
Gastric (STAD)   & 368 & chromosomal-instability vs genomically-stable/MSI \\
Lung adeno (LUAD)& 402 & amp drivers (EGFR/MET/MYC) vs whole-genome-doubling \\
Esophageal (ESCA)& 124 & RTK amp (ERBB2/CCNE1) vs arm-level aneuploidy \\
\bottomrule
\end{tabular}
\end{center}

\section{Results}

\subsection{The decomposition reproduces and recovers real biology}
In 6 of 8 cohorts the GSVD produced the two antisymmetric arraylets the theory
predicts ($c_{\text{first}}\!\approx\!1$, $c_{\text{last}}\!\approx\!0$) with the
two patient-score vectors close to orthogonal (cosine $0.002$--$0.089$). The
top-loading genomic regions of the two arraylets fell onto the known
cancer-specific CN axes, and the method also correctly isolated the demographic
(sex / X-chromosome) pattern as a separate, non-prognostic axis. \textbf{The GSVD
machinery and the ``two orthogonal CN programs'' framing are real and general} ---
the decomposition does what the paper says it does. This is the durable
contribution of the work.

\subsection{The prognostic strength does not generalize}
Table~\ref{tab:main} reports the survival results across all eight cohorts.

\begin{table}[h]
\centering
\small
\caption{Survival concordance of the GSVD arraylets versus baselines, eight
cohorts. $\C=0.5$ is chance. $N_m$ = matched patients; $N_s$ = patients with
survival data; $\C_1$/$\C_2$ = first/second arraylet; $\C_{\text{comb}}$ =
combined; $\cos_\perp$ = cross-arraylet orthogonality; age and Cox-20 are the
baselines the predictor must beat.}
\label{tab:main}
\begin{tabular}{@{}lrrrrrrrrcc@{}}
\toprule
Cohort & $N_m$ & $N_s$ & $\C_1$ & $\C_2$ & $\C_{\text{comb}}$ & $\cos_\perp$ & age $\C$ & Cox-20 & $>$age & $>$Cox \\
\midrule
GLIOMA & 976 & 899 & 0.275 & 0.510 & 0.277 & 0.002 & 0.744 & 0.737 & N & N \\
BRCA   & 978 & 920 & 0.542 & 0.462 & 0.542 & 0.049 & 0.614 & 0.644 & N & N \\
OV     & 441 & 391 & 0.561 & 0.491 & 0.557 & 0.038 & 0.608 & 0.636 & N & N \\
SARC   & 232 & 196 & 0.496 & 0.490 & 0.498 & 0.377 & 0.611 & 0.707 & N & N \\
UCEC   & 497 & 456 & 0.333 & 0.505 & 0.333 & 0.004 & 0.603 & 0.708 & N & N \\
STAD   & 368 & 345 & 0.482 & 0.506 & 0.482 & 0.007 & 0.554 & 0.592 & N & N \\
LUAD   & 402 & 345 & 0.469 & 0.464 & 0.469 & 0.089 & 0.526 & 0.644 & N & N \\
ESCA   & 124 & 106 & 0.503 & 0.507 & 0.502 & 0.071 & 0.628 & 0.714 & N & N \\
\bottomrule
\end{tabular}
\end{table}

Across all eight cohorts --- including ovarian, the adult cancer whose biology
most resembles neuroblastoma's (copy-number-driven, near-universal TP53 loss,
minimal point-mutation drivers) --- the tumor-exclusive arraylet's survival
concordance sits at chance (range $0.275$--$0.561$, median $\approx0.50$), and
\textbf{age and a plain supervised Cox model beat it in every single cohort
($0/8$ and $0/8$)}. The paper's headline $\C\approx0.77$--$0.80$, beating
standard of care, did not reproduce in any other CN-driven cancer.

\section{Interpretation: good science, cleanly stated}
\label{sec:interp}

This is a textbook, constructive replication outcome that separates into two
layers.

\textbf{(1) The method is real and the decomposition is a durable contribution.}
The two-orthogonal-mechanism GSVD reproduces across cancers and lands on the right
biology. As a descriptive, interpretable tool for extracting independent
copy-number programs --- and for cleanly separating biological signal from
demographic artifact --- it does what is claimed.

\textbf{(2) The strong prognostic claim is cohort-specific, not universal.} The
$\C\approx0.77$--$0.80$ result is consistent with two ordinary, non-pejorative
explanations, which are not mutually exclusive:
\begin{itemize}[leftmargin=1.4em]
  \item \emph{Neuroblastoma-specific biology.} NBL is the textbook
  copy-number-driven tumor; its prognosis is genuinely written into the CN profile
  (MYCN/2p24.3 --- which $u_{1,1}$ maps to directly --- plus segmental 1p/11q/17q
  and ID2, which $u_{1,101}$ maps to). Most adult cancers carry prognostic
  information in mutation and methylation layers that a CN-only arraylet cannot
  see. The method may work in NBL precisely because copy number \emph{is} the
  disease there.
  \item \emph{In-sample optimism.} The reported concordance is measured on the
  same discovery patients from which the predictor was derived, without
  cross-validation and without a supervised baseline for comparison. In-sample
  concordance on a fitted predictor is expected to exceed its out-of-sample value;
  this is the universal reason a striking single-cohort predictor requires external
  replication.
\end{itemize}
Neither explanation implies anything beyond the need for independent validation
--- which is precisely what replication provides.

\section{Limitations}
\label{sec:limits}
Two limitations bound how hard our open-data test could push the original claim,
and we state them plainly.

\textbf{(a) Feature resolution.} The paper uses $\sim$2.83 million one-kilobase
bins; our open substrate is GDC ``Masked Copy Number Segment'' (DNAcopy CBS),
which we tiled at one-megabase resolution ($\sim$2{,}633 autosomal bins) ---
roughly three orders of magnitude coarser. Because the segment data are
piecewise-constant, finer binning would mostly re-copy segment values rather than
add true within-segment resolution; nonetheless, focal driver amplifications
smaller than a called segment (e.g. MYCN, CCNE1, ERBB2, MDM2/CDK4) can be smeared
across a broad bin and their signal averaged down. Consequently our test is a fair
probe of arm- and broad-scale copy-number prognosis but a \emph{weaker} probe of
the focal-amplification prognosis the original may depend on. This does not change
the verdict --- the decomposition still reproduces and the supervised baseline
still wins --- but it is an honest limit on how stringently we tested the
focal-driver claim, and it is a further reason the decisive experiment is the
native 1-kb WGS, evaluated out-of-sample.

\textbf{(b) Unsupervised by construction.} The arraylets are chosen by the
decomposition's own variance/orthogonality criteria and never see the survival
labels; the supervised penalized-Cox baseline does. An unsupervised
maximum-tumor-contrast axis predicts outcome only when the dominant aberration
axis happens to coincide with the prognostic axis --- a coincidence that can hold
in copy-number-driven disease and break elsewhere. This is a property of the
method's design, not a defect, but it explains why the supervised baseline wins in
every cohort.

\textbf{(c) Resolution is intrinsic to the open source, not a tunable knob.} The
SNP6 ``Masked Copy Number Segment'' product is natively segment-grained; tiling at
one megabase is the appropriate choice \emph{for array-segment data}. The
$\sim$1000$\times$ gap to the paper's 1-kb WGS is therefore inherent to the open
data source --- we cannot present the method with the focal sub-megabase signal on
open data at all. This is precisely why the decisive experiment requires the native
1-kb WGS (Route A / dbGaP).

\textbf{(d) Matched-normal modality.} The blood layer is the blood-derived-normal
aliquot's own SNP6 segment CN --- the correct \emph{construction} (a genuine
independent germline genome, $Z_2$), but a different measurement modality
(array-segment CN) than the paper's WGS read-depth $Z_2$.

\textbf{(e) Out-of-sample by cohort, not by the original split; artifact hygiene.}
Our concordances are out-of-sample in the sense of being independent cohorts; only
a rerun on the authors' own patients with cross-validation gives the clean
in-sample-vs-out-of-sample comparison to their 0.80. (Minor: the per-cohort result
files carry a templated \texttt{cohort:"TCGA-BRCA"} label from the build scaffold;
the filename is authoritative and the numbers are unaffected.)

\textbf{The honest boundary.} The prognostic null established here is narrow:
\emph{on segment-resolution open array data, these predictors do not beat baselines
out of sample}. It does \emph{not} establish that the method fails on native 1-kb
WGS in copy-number-driven disease. Crucially, the geometry and two-mechanism
reproductions are robust to all of caveats (a)--(e) --- they do not depend on
resolution or modality --- which is exactly why they hold while the prognostic
layer does not.

\section{Why the prognostic layer failed --- a precise diagnosis}
\label{sec:whyfail}
The survival null is not vague ``the method doesn't work'' --- it is specific and
overdetermined. Four independent causes each push toward chance-level prognosis;
together they guarantee it.

\begin{enumerate}[leftmargin=1.6em]
  \item \textbf{The geometry's success is exactly what allows the prognosis to
  fail.} The tumor-exclusive arraylet ($c_{\text{last}}\!\to\!0$) is by
  construction the copy-number pattern of \emph{maximum tumor-vs-normal contrast}
  --- a statement about the \emph{magnitude} of somatic aberration, not about
  \emph{which} aberrations drive outcome. The biggest contrast is dominated by
  \emph{passenger} CN changes (broad arm-level gains/losses every tumor
  accumulates), not the handful of \emph{driver} events that predict survival. So
  the method reliably extracts a real axis that is prognostically diffuse.
  \item \textbf{It is unsupervised, so nothing ever pointed it at survival.} The
  arraylets are chosen by the decomposition's own variance/orthogonality criteria
  and never see outcome labels. The penalized-Cox baseline does see the labels ---
  which is precisely why it wins in every cohort. An unsupervised axis predicts
  outcome only when the dominant aberration axis happens to coincide with the
  prognostic axis: a coincidence that can hold in copy-number-driven disease and
  breaks elsewhere.
  \item \textbf{Resolution.} The paper uses $\sim$2.83M one-kilobase features; our
  open substrate is one-megabase ($\sim$2{,}633 bins), $\sim$1000$\times$ coarser.
  The segment data are piecewise-constant, so finer binning would mostly re-copy
  segment values --- but focal driver amplifications smaller than a called segment
  (MYCN, CCNE1, ERBB2, MDM2/CDK4) can be smeared into a broad lukewarm bin and
  averaged out. This is the single biggest reason the open-data test cannot fairly
  evaluate the focal-driver form of the claim.
  \item \textbf{Wrong biology in most cohorts.} Adult glioma is driven by point
  mutations and methylation (IDH, 1p/19q, MGMT) that copy number of \emph{any}
  resolution cannot encode --- hence $\C=0.28$ there (an anti-correlated noise
  axis, worse than chance).
\end{enumerate}

\textbf{The diagnostic tell that the method is not broken:} the geometry
reproduced flawlessly ($c_{\text{first}}\approx1$, $c_{\text{last}}\ll1$,
orthogonality $\approx0$ in most cohorts). The linear algebra does exactly what it
claims. What did not transfer is the \emph{leap} from ``this axis is the largest
tumor-specific contrast'' to ``this axis predicts survival'' --- a leap that holds
only under the original's specific conditions (focal 1-kb resolution +
copy-number-driven NBL biology + likely in-sample evaluation). That is a
conditional, well-understood failure, not a refutation of the decomposition.

\section{What would settle it}
\begin{itemize}[leftmargin=1.4em]
  \item \textbf{Run the authors' pipeline on the authors' NBL data with
  cross-validation and a supervised-Cox baseline added.} This requires Orly
  Alter's processed Datasets~1--3 (offered on request) and/or a direct dbGaP DAR
  for phs000467 (an NIH-funded PI is independently eligible). This single
  experiment distinguishes ``real but NBL-specific'' from ``in-sample optimism.''
  \item \textbf{Pediatric analog (optional).} Medulloblastoma (Group~3 MYC vs
  Group~4 isodicentric-17q) is the closest pediatric CN-driven two-mechanism
  cohort; testing it would sharpen ``pediatric-CN-specific'' versus
  ``NBL-specific.'' Requires non-GDC data sourcing (germline CN often controlled).
\end{itemize}

\section{Summary}
We independently reproduced the Alter multitensor GSVD on its native data type
(matched tumor and blood genome-wide copy number) across eight cancers, using a
from-scratch engine and open data, with zero controlled-access dependencies. The
decomposition and its two-orthogonal-mechanism structure reproduce robustly and
recover real, cancer-specific copy-number biology --- a genuine methodological
contribution. The strong prognostic claim (concordance $\approx0.77$--$0.80$,
beating standard of care) did not generalize: in every cohort tested, including
the copy-number-driven cancers most analogous to neuroblastoma, the predictor
performed at chance and was outperformed by patient age and an ordinary supervised
model. The most likely reading is that the prognostic result is specific to
neuroblastoma's copy-number-driven biology and/or reflects in-sample optimism in a
single discovery cohort --- both of which call for the standard remedy:
replication on the original data with cross-validation and a supervised baseline.
That experiment is well-defined and within reach. This is replication working as
intended: separating the durable methodological contribution from the
cohort-bound empirical claim, and pointing to the one decisive next experiment.

\vspace{1em}
\hrule
\vspace{0.5em}
{\footnotesize
\textbf{Compute \& data.} All analyses on uicgpu (8$\times$A100, free).
Data: NCI Genomic Data Commons open tier (``Masked Copy Number Segment,''
DNAcopy/SNP6) for eight TCGA projects; survival metadata from the GDC clinical
endpoint. Engine and results:
\texttt{/data/stevens/alter-pathB/} (code: \texttt{gsvd\_reference.py},
\texttt{survival\_stats.py}, \texttt{f4--f6\_*.py}; results:
\texttt{f4\_glioma\_results.json}, \texttt{f5\_brca\_results.json},
\texttt{f6\_\{OV,SARC,UCEC,STAD,LUAD,ESCA\}\_results.json}). No dbGaP /
controlled-access data were used.}

\end{document}
